%%%%%%%%%%%%%%%%%%%%%%%%%%%%%%%%%%%%%%%%%%%%%%%%%%%%%%%%%%%%%%%%%%%%%%%%%%%%%%%%
\chapter{Type Safety}\label{ch:type-safety}
%%%%%%%%%%%%%%%%%%%%%%%%%%%%%%%%%%%%%%%%%%%%%%%%%%%%%%%%%%%%%%%%%%%%%%%%%%%%%%%%

\change[inline]{Add fluff text between lemmata.}

\section{Preservation}

\begin{lemma}[\rocqicon~Preservation for $\equiv$] \label{lemma:preservation-for-struct-cong}
  If $P \equiv Q$, then $\Gamma \vdash P$ if and only if $\Gamma \vdash Q$.
\end{lemma}
\begin{proof}
  By induction on the derivation of $P \equiv Q$.
\end{proof}

\begin{lemma}[\rocqicon~Partitions for well-formed substitutions] \label{lemma:subst-partition}
  If $\Gamma = \Gamma_1 \circ \Gamma_2$ and $\Gamma \vdash \sigma : \Delta$, then
  there exist $\Delta_1$ and $\Delta_2$ such that
  \begin{enumerate}[label=(\roman*)]
    \item $\Delta = \Delta_1 \circ \Delta_2$ and
    \item $\Gamma_1 \vdash \sigma : \Delta_1$ and
    \item $\Gamma_2 \vdash \sigma : \Delta_2$.
  \end{enumerate}
\end{lemma}
\begin{proof}
  By induction on the length of the tail of $\Gamma$.
\end{proof}

\begin{lemma}[\rocqicon~Substitution]\label{lemma:substitution}\mbox{}
  \begin{enumerate}[label=(\roman*)]
    \item If $\Gamma \vdash P$, then, for all substitutions $\sigma$ such that $\Gamma \vdash \sigma : \Delta$,
          it follows $\Delta \vdash P\sigma$.
    \item If $\Gamma \vdash M : A$, then, for all substitutions $\sigma$ such that $\Gamma \vdash \sigma : \Delta$,
          it follows $\Delta \vdash M\sigma : A$.
    \item If $\Gamma \vdash S : A$, then, for all substitutions $\sigma$ such that $\Gamma \vdash \sigma : \Delta$,
          it follows $\Delta \vdash S\sigma : A$.
  \end{enumerate}
\end{lemma}
\begin{proof}
  By induction on the typing derivation. We will only highlight the case
  for \textsc{t-cut}, the remaining cases follow immediately or analogously.

  We may assume that the names of free and bound variables are unique, i.e.
  $\parcomp{x}{y}{P}{Q}\sigma = \parcomp{x}{y}{P\sigma}{Q\sigma}$.
  Given
  \begin{center}
    \begin{prooftree}
      \hypo{\Gamma_1, x : A \vdash P}
      \hypo{\Gamma_2, y : A^\bot \vdash Q}
      \infer2[\textsc{t-cut}]{\Gamma \vdash \parcomp{x}{y}{P}{Q}}
    \end{prooftree}
  \end{center}
  with $\Gamma = \Gamma_1 \circ \Gamma_2$. By induction hypothesis,
  we have for all $\Delta$ and $\sigma$, if $\Gamma_1, x : A \vdash \sigma : \Delta$,
  then $\Delta \vdash P\sigma$, and similarly for $Q$.
  Using Lemma \ref{lemma:subst-partition}, there exist $\Delta_1$ and $\Delta_2$ such that
  \[
    \Gamma_1 \vdash \sigma : \Delta_1 \quad \text{and} \quad
    \Gamma_2 \vdash \sigma : \Delta_2 \quad \text{and} \quad
    \Delta = \Delta_1 \circ \Delta_2.
  \]
  Linear typing further yields
  \[
    \Gamma_1, x : A \vdash (x \coloneq x) \cdot \sigma : \Delta_1, x : A
  \]
  and
  \[
    \Gamma_2, y : A^\bot \vdash (y \coloneq y) \cdot \sigma : \Delta_2, y : A^\bot
  \]
  Thus, using the induction hypothesis,
  \begin{center}
    \begin{prooftree}
      \hypo{\Delta_1, x : A \vdash P\sigma}
      \hypo{\Delta_2, y : A^\bot \vdash Q\sigma}
      \infer2[\textsc{t-cut}]{\Delta \vdash \parcomp{x}{y}{P\sigma}{Q\sigma}}
    \end{prooftree}
  \end{center}
  where we use the fact that $P\sigma = P((x \coloneq x) \cdot \sigma)$.
\end{proof}

\begin{theorem}[\rocqicon~Preservation]
  If $\Gamma \vdash P$ and $P \contr Q$, then $\Gamma \vdash Q$.
\end{theorem}
\begin{proof}
  By induction on the derivation of the premise $P \contr Q$.
  Most cases follow immediately by inversion on the typing relation.
  The interesting cases are \textsc{r-axcut}, \textsc{r-seq} and \textsc{r-struct}.

  \begin{enumerate}[itemsep=1em, label=Case:, leftmargin=*]
    \item (\textsc{r-axcut}) \\
      We have
      \[
        P = \parcomp{x}{y}{\link{x}{M}}{R} \contr R[M/y] = Q.
      \]
      By inversion on the typing derivation
      \[
        \Gamma \vdash \parcomp{x}{y}{\link{x}{M}}{R}
      \]
      we obtain
      \[
        x : A \vdash x : A \quad \text{and} \quad
        \Delta_1 \vdash M : A^\bot \quad \text{and} \quad
        \Delta_2, y : A^\bot \vdash R
      \]
      with $\Gamma = \Delta_1 \circ \Delta_2$. Defining
      \[
        \sigma \coloneq (y \coloneq M) \cdot id_{\Delta_2},
      \]
      we conclude $R[M/y] = R\sigma$ using Lemma \ref{lemma:strict-domains}.
      Further, we can derive
      \[
        \Delta_2, y : A^\bot \vdash \sigma : \Gamma.
      \]
      by the disjointness of $\Delta_1$ and $\Delta_2$.
      Thus, using Lemma \ref{lemma:substitution}, we obtain $\Gamma \vdash R \sigma$.
    \item (\textsc{r-seq}) \\
      Analogous to the previous case.
    \item (\textsc{r-struct}) \\
      The derivation ends with the following shape:
      \begin{center}
        \begin{prooftree}
          \hypo{P \equiv P'}
          \hypo{P' \contr Q'}
          \hypo{Q' \equiv Q}
          \infer3[]{P \contr Q}
        \end{prooftree}
      \end{center}
      By induction hypothesis, we have, for any $\Gamma$,
      if $\Gamma \vdash P'$, then $\Gamma \vdash Q'$. Using Lemma
      \ref{lemma:preservation-for-struct-cong}, we obtain
      $\Gamma \vdash P'$. Thus, by IH, $\Gamma \vdash Q'$. Finally,
      using Lemma \ref{lemma:preservation-for-struct-cong} again, we
      derive the desired conclusion $\Gamma \vdash Q$.
  \end{enumerate}
\end{proof}

\section{Terminated Processes}
\subsection{Canonical Cut Forms and Link Lists}
\begin{definition}[Canonical cut form]
  The set of cuts that is in \emph{canonical cut form} is inductively defined by the following clauses.
  \begin{enumerate}[label=(\roman*)]
    \item If $L$ and $R$ are not cuts, then
          $\parcomp{x}{y}{L}{R}$ is in canonical cut form.
    \item If $L$ is not a cut and $P$ is in canonical cut form,
          then $\parcomp{x}{y}{L}{P}$ is also in canonical cut form.
  \end{enumerate}
\end{definition}

\begin{remark}
  A cut that is in canonical cut form is a degenerate tree of $\textsc{cut}$s
  with link processes or sequences as leaves. That is why we sometimes refer
  to a cut in canonical cut form as a cut list.
\end{remark}

\begin{lemma}[\rocqicon~Appending cut lists] \label{lemma:cut-list-append}
  If $\Gamma \vdash \parcomp{x}{y}{P}{Q}$ and both $P$ and $Q$ are in canonical
  cut form, then there exists an $R$ in canonical cut form such that $\parcomp{x}{y}{P}{Q} \equiv R$.
\end{lemma}
\begin{proof}
  By induction on the depth of $P$ and appropriate application of structural congruences for cuts.
\end{proof}

\begin{lemma}[\rocqicon~Existence of canonical cut forms] \label{lemma:canonical-cut-forms}
  If $\Gamma \vdash \parcomp{x}{y}{L}{R}$, then there exists
  a process $P$ such that $\parcomp{x}{y}{L}{R} \equiv P$ and $P$
  is in canonical cut form.
\end{lemma}
\begin{proof}
  By induction on the derivation of $\Gamma \vdash \parcomp{x}{y}{L}{R}$, followed
  by a case analysis on the structure of $L$ and $R$. Excluding the vacuous
  cases which are impossible due to typing, we end up with five valid cases:
  four where one premise is a link or sequencing and the other is a cut,
  and a fifth, where both $L$ and $R$ are cuts.

  \begin{enumerate}[itemsep=1em, label=Case:, leftmargin=*]
  \item ($L = \link{M_1}{M_2}$ and $R = \parcomp{a}{b}{Q_1}{Q_2}$) \\
    By induction hypothesis, there exists an $R'$ such that $R \equiv R'$ and
    $R'$ is in canonical cut form. Thus,
    \[
      \parcomp{x}{y}{L}{R} \equiv \parcomp{x}{y}{\link{M_1}{M_2}}{R'} \eqcolon P
    \]
    where $P$ is in canonical cut form by definition.
    The other three cases in which exactly one of $L$ or $R$
    is a cut follow analogously.

  \item ($L = \parcomp{a}{b}{L_1}{L_2}$ and $R = \parcomp{c}{d}{R_1}{R_2}$) \\
    By induction hypothesis, there exists $L'$ and $R'$ such that
    \[
      L \equiv L' \quad \text{and} \quad R \equiv R',
    \]
    and both $L'$ and $R'$ are in canonical cut form.
    Thus, $\parcomp{x}{y}{L}{R} \equiv \parcomp{x}{y}{L'}{R'}$. The result then
    follows by applying Lemma \ref{lemma:cut-list-append}.
  \end{enumerate}
\end{proof}

A special case of a cut list is a \emph{link list}, which,
informally, is a degenerate tree of cuts whose leaves are link processes
of the form $\link{x_i}{M_i}$ (or symmetrically $\link{M_i}{x_i}$).

\begin{definition}[Link list]
  The set of \emph{link lists} and their associated \emph{spines} (sequences of names) are
  inductively defined by the following clauses.
  \begin{enumerate}[label=(\roman*)]
    \item If $P_1 \equiv \link{x_1}{M_1}$ and $P_2 \equiv \link{x_2}{M_2}$ such that
          $M_1$ and $M_2$ are not futures, then $\parcomp{x}{y}{P_1}{P_2}$
          is a link list with spine $spine(P) = [x_1, x_2]$.
    \item If $L \equiv \link{a}{M}$ such that $M$ is not a future and
          $Q$ is a link list with spine $spine(Q) = \vec{x}$, then 
          $P = \parcomp{a}{b}{L}{Q}$ is a link list with spine $spine(P) = x :: \vec{x}$.
  \end{enumerate}
\end{definition}

We can think of a link list as a pool of threads. If every channel
in the spine is a free occurrence, then every thread is waiting on external communication,
and we consider the process terminated. Otherwise, if at least one channel
in the spine is bound, it can pass its message to another thread in the pool.

\begin{lemma}[\rocqicon~Reducibility of link lists] \label{lemma:link-list-reducibility}
  Let $\Gamma \vdash P$ be a well-typed link list. If some channel
  $a \in spine(P)$ is bound in $P$, then $P \contr P'$ for some process $P'$.
\end{lemma}
\begin{proof}[Proof Sketch]
  The proof proceeds by strong induction on the length of the spine. It relies
  on several auxiliary lemmas regarding the reducibility of link lists with specific
  spine properties. Due to their technical nature, we omit these intermediate
  results here for brevity, referring the interested reader to the formalization.

  However, the general intuition is as follows: If $a \in spine(P)$ occurs
  bound in $P$, then, without loss of generality, we may assume that $a$ is bound 
  by the leftmost scope; i.e., we are looking at the following situation:
  \[
    (\nu x a)(L \mid E[\parcomp{c}{d}{\link{a}{M}}{P}]).
  \]
  Here, the context $E$ describes a link list prefix whose tail is instantiated 
  with the link list $\parcomp{c}{d}{\link{a}{M}}{P}$. Through a series of
  structural congruences, it is possible to commute the binder $(\nu x a)$ 
  inwards to the link $\link{a}{M}$ to eventually yield a process
  \[
    E'[\parcomp{a}{x}{\link{a}{M}}{P'}],
  \]
  for a new list prefix $E'$ and tail $P'$. This process may now reduce via
  \textsc{r-axcut}.
\end{proof}

\subsection{Final Processes}

\begin{definition}[Final process]
  A process $P$ is \emph{final} if it satisfies one of the following conditions:
  \begin{enumerate}[label=(\roman*)]
    \item $P = \Lambda$.
    \item $P$ is of the form $\link{x}{M}$ or $\link{M}{x}$.
    \item There exists a link list $Q$ such that $P \equiv Q$ and
          every channel $a \in spine(Q)$ is free in $Q$.
  \end{enumerate}
\end{definition}

\begin{lemma}[\rocqicon~Final processes are irrecudible]
  If $\Gamma \vdash P$, then $P$ is final if and only if
  $P$ is irreducible.
\end{lemma}

\change[inline]{Give high-level explanation of the proof;
explain difficulty (that is, no member of the same equivalence can contain a redex)}

\section{Progress}

\begin{lemma}[\rocqicon~Reducibility of canonical cut forms] \label{lemma:cut-list-reducibility}
  If $\Gamma \vdash P$ is a well-typed cut list,
  then $P \contr P'$ for some process $P'$ or $P$ is a link list.
\end{lemma}
\begin{proof}[Proof Sketch]
  The proof is by induction on the structure of the cut list $P$.
  We outline the underlying mechanism:
  \begin{enumerate}
    \item As long as a cut list contains a leaf that is not of the form $\link{x}{M}$ (or $\link{M}{x}$),
          well-typedness guarantees that this leaf can step via \textsc{r-seq} or
          a $\beta$-reduction.

    \item If all leaves have the shape $\link{x}{M}$ (or $\link{M}{x}$), then $P$ is a link
          list by definition, which completes the proof.
  \end{enumerate}
\end{proof}

\begin{corollary}[\rocqicon~Progress for canonical cut forms] \label{cor:progress-for-cut-list}
  If $\Gamma \vdash P$ is a cut list, then $P$ is final
  or $P \contr P'$ for some process $P'$.
\end{corollary}
\begin{proof}
  Using Lemma \ref{lemma:cut-list-reducibility}, we consider the following two cases:
  \begin{enumerate}[label=(\roman*)]
    \item $P \contr P'$ for some $P'$. The result follows immediately.
    \item $P$ is a link list. If all names in its spine are free occurrences,
          then $P$ is final by definition. Otherwise, some channel in $spine(P)$
          is bound in $P$, so $P \contr P'$ for some $P'$
          by Lemma \ref{lemma:link-list-reducibility}.
  \end{enumerate}
\end{proof}

\begin{theorem}[\rocqicon~Progress]
  If $\Gamma \vdash P$, then $P \contr P'$ for some process $P'$ or
  $P$ is final.
\end{theorem}
\begin{proof}
  By inversion on the typing derivation of $\Gamma \vdash P$.
  \begin{enumerate}[itemsep=1em, label=Case:, leftmargin=*]
    \item (\textsc{t-ax}) \\
      We have $P = \link{M_1}{M_2}$ and observe, that a well-typed link process
      may either $\beta$-reduce (yielding $P'$) or is of the form $\link{x}{M}$ (or $\link{M}{x}$),
      in which case $P$ is final.

    \item (\textsc{t-cut}) \\
      We have $P = \parcomp{x}{y}{P_1}{P_2}$. Using Lemma
      \ref{lemma:canonical-cut-forms}, there exists a canonical cut form
      $P'$ such that $P \equiv P'$. By Lemma \ref{lemma:preservation-for-struct-cong},
      it follows $\Gamma \vdash P'$.
      Now, applying Lemma \ref{cor:progress-for-cut-list}, we consider
      the following two subcases for $P'$:
      \begin{enumerate}[label=(\roman*)]
        \item $P'$ is final. Since $P$ is a cut, it follows by inversion
              on $P \equiv P'$, that $P'$ is a link list with no bound
              name in its spine. Thus, by definition, $P$ is final.
        \item $P' \contr R$ for some $R$. Consequently, $P \contr R$ via \textsc{r-struct}.
      \end{enumerate}

    \item (\textsc{t-seq}) \\
      Now, $P = Q \seq xM$, which reduces to $P \contr Q[M/x] \eqcolon P'$.

    \item (\textsc{t-nil}) \\
      We have $P = \Lambda$, which is already final.
  \end{enumerate}
\end{proof}

\begin{corollary}[\rocqicon~Progress for closed processes]
  If $\vdash P$, then $P \contr P'$ for some process $P'$ or
  $P = \Lambda$.
\end{corollary}
\begin{proof}
  Immediate from Lemma \ref{lemma:strict-domains} and the observation that
  $\Lambda$ is the only closed, final process.
\end{proof}
